\subsection{\label{sec.gf.nips}Non-integer powers}

\subsubsection{Definition}

Now, let us again recall Example 2 from Section \ref{sec.gf.exas}. In order to
fully justify that example, we still need to explain what $\sqrt{1-4x}$ is.

More generally, let us try to define non-integer powers of FPSs (since square
roots are just $1/2$-th powers). Thus, we are trying to solve the following problem:

\begin{statement}
\textit{Problem:} Devise a reasonable definition of the $c$-th power $f^{c}$
for any FPS $f\in K\left[  \left[  x\right]  \right]  $ and any $c\in K$.
\end{statement}

Here, \textquotedblleft reasonable\textquotedblright\ means that it should
have some of the properties we would expect:

\begin{itemize}
\item It should not conflict with the existing notion of $f^{c}$ for
$c\in\mathbb{N}$. That is, if $c\in\mathbb{N}$, then our new definition of
$f^{c}$ should yield the same result as the existing meaning that $f^{c}$ has
in this case (namely, $\underbrace{ff\cdots f}_{c\text{ times}}$). The same
should hold for $c\in\mathbb{Z}$ when $f$ is invertible.

\item Rules of exponents should hold: i.e., we should have%
\begin{equation}
f^{a+b}=f^{a}f^{b},\ \ \ \ \ \ \ \ \ \ \left(  fg\right)  ^{a}=f^{a}%
g^{a},\ \ \ \ \ \ \ \ \ \ \left(  f^{a}\right)  ^{b}=f^{ab}
\label{eq.sec.gf.nips.rules-of-exps}%
\end{equation}
for all $a,b\in K$ and $f,g\in K\left[  \left[  x\right]  \right]  $.

\item For any positive integer $n$ and any FPS $f\in K\left[  \left[
x\right]  \right]  $, the $1/n$-th power $f^{1/n}$ should be an $n$-th root of
$f$ (that is, an FPS whose $n$-th power is $f$). (This actually follows from
the previous two properties, since we can apply the rule $\left(
f^{a}\right)  ^{b}=f^{ab}$ to $a=1/n$ and $b=n$.)
\end{itemize}

Clearly, we cannot solve the above problem in full generality:

\begin{itemize}
\item The power $0^{-1}$ cannot be reasonably defined (unless $K$ is trivial).
Indeed, $0^{-1}\cdot0^{1}$ would have to equal $0^{-1+1}=0^{0}=1$, but this
would contradict $0^{-1}\cdot0^{1}=0^{-1}\cdot0=0$.

\item The power $x^{1/2}$ cannot be reasonably defined either (unless $K$ is
trivial). Indeed, there is no FPS whose square is $x$. This will be proved in
Exercise \ref{exe.fps.not-squares} \textbf{(a)}.

\item Even the power $\left(  -1\right)  ^{1/2}$ cannot always be defined:
There is no guarantee that $K$ contains a square root of $-1$ (and if $K$ does
not, then it is easy to see that $K\left[  \left[  x\right]  \right]  $ does neither).
\end{itemize}

However, all we want is to make sense of $\sqrt{1-4x}$, so let us restrict
ourselves to FPSs whose constant term is $1$. Using the notation from
Definition \ref{def.fps.Exp-Log-maps} \textbf{(b)}, we are thus moving on to
the following problem:

\begin{statement}
\textit{More realistic problem:} Devise a reasonable definition of the $c$-th
power $f^{c}$ for any FPS $f\in K\left[  \left[  x\right]  \right]  _{1}$ and
any $c\in K$.
\end{statement}

Besides imposing the above wishlist of properties, we want this $c$-th power
$f^{c}$ itself to belong to $K\left[  \left[  x\right]  \right]  _{1}$, since
otherwise the iterated power $\left(  f^{a}\right)  ^{b}$ in our rules of
exponents might be undefined.

It turns out that this is still too much to ask. Indeed, if $K=\mathbb{Z}/2$,
then the FPS $1+x\in K\left[  \left[  x\right]  \right]  _{1}$ has no square
root (you get to prove this in Exercise \ref{exe.fps.not-squares}
\textbf{(c)}), so its $1/2$-th power $\left(  1+x\right)  ^{1/2}$ cannot be
reasonably defined.

However, if we assume (as in Convention \ref{conv.fps.exp.K-Q-alg}) that $K$
is a commutative $\mathbb{Q}$-algebra, then we get lucky: Our
\textquotedblleft more realistic problem\textquotedblright\ can be solved in
(at least) two ways: \medskip

\textit{1st solution:} We define
\[
\left(  1+x\right)  ^{c}:=\sum_{k\in\mathbb{N}}\dbinom{c}{k}x^{k}%
\ \ \ \ \ \ \ \ \ \ \text{for each }c\in K,
\]
in order to make Newton's binomial formula (Theorem \ref{thm.fps.newton-binom}%
) hold for arbitrary exponents\footnote{Note that $\dbinom{c}{k}%
=\dfrac{c\left(  c-1\right)  \left(  c-2\right)  \cdots\left(  c-k+1\right)
}{k!}$ is well-defined since $K$ is a commutative $\mathbb{Q}$-algebra.}.
Subsequently, we define%
\begin{equation}
f^{c}:=\left(  1+x\right)  ^{c}\left[  f-1\right]
\ \ \ \ \ \ \ \ \ \ \text{for any }f\in K\left[  \left[  x\right]  \right]
_{1}\text{ and }c\in K \label{eq.sec.gf.nips.newton-def2}%
\end{equation}
(in order to have $\left(  1+g\right)  ^{c}=\sum_{k\in\mathbb{N}}\dbinom{c}%
{k}g^{k}$ hold not only for $g=x$, but also for all $g\in K\left[  \left[
x\right]  \right]  _{0}$).

It is clear that the FPS $f^{c}$ is well-defined in this way. However, proving
that this definition satisfies all our wishlist (particularly the rules of
exponents (\ref{eq.sec.gf.nips.rules-of-exps})) is highly nontrivial. Some of
this is done in \cite[\S 7.12]{Loehr-BC}, but it is still a lot of work.

Thus, we shall discard this definition of $f^{c}$, and instead take a
different way: \medskip

\textit{2nd solution:} Recall the mutually inverse group isomorphisms%
\begin{align*}
\operatorname*{Exp}  &  :\left(  K\left[  \left[  x\right]  \right]
_{0},+,0\right)  \rightarrow\left(  K\left[  \left[  x\right]  \right]
_{1},\cdot,1\right)  \ \ \ \ \ \ \ \ \ \ \text{and}\\
\operatorname*{Log}  &  :\left(  K\left[  \left[  x\right]  \right]
_{1},\cdot,1\right)  \rightarrow\left(  K\left[  \left[  x\right]  \right]
_{0},+,0\right)
\end{align*}
from Theorem \ref{thm.fps.Exp-Log-group-iso}. Thus, for any $f\in K\left[
\left[  x\right]  \right]  _{1}$ and any $c\in\mathbb{Z}$, the equation%
\[
\operatorname*{Log}\left(  f^{c}\right)  =c\operatorname*{Log}f
\]
holds (since $\operatorname*{Log}$ is a group homomorphism). This suggests
that we define $f^{c}$ for all $c\in K$ by the same equation. In other words,
we define $f^{c}$ for all $c\in K$ by setting $f^{c}=\operatorname*{Exp}%
\left(  c\operatorname*{Log}f\right)  $ (since the map $\operatorname*{Exp}$
is inverse to $\operatorname*{Log}$). And this is what we shall do now:

\begin{definition}
\label{def.fps.power-c}Assume that $K$ is a commutative $\mathbb{Q}$-algebra.
Let $f\in K\left[  \left[  x\right]  \right]  _{1}$ and $c\in K$. Then, we
define an FPS%
\[
f^{c}:=\operatorname*{Exp}\left(  c\operatorname*{Log}f\right)  \in K\left[
\left[  x\right]  \right]  _{1}.
\]

\end{definition}

This definition of $f^{c}$ does not conflict with our original definition of
$f^{c}$ when $c\in\mathbb{Z}$ because (as we said) the original definition of
$f^{c}$ already satisfies $\operatorname*{Log}\left(  f^{c}\right)
=c\operatorname*{Log}f$ and therefore $f^{c}=\operatorname*{Exp}\left(
c\operatorname*{Log}f\right)  $.

Moreover, Definition \ref{def.fps.power-c} makes the rules of exponents hold:

\begin{theorem}
\label{thm.fps.power-c.rules}Assume that $K$ is a commutative $\mathbb{Q}%
$-algebra. For any $a,b\in K$ and $f,g\in K\left[  \left[  x\right]  \right]
_{1}$, we have
\[
f^{a+b}=f^{a}f^{b},\ \ \ \ \ \ \ \ \ \ \left(  fg\right)  ^{a}=f^{a}%
g^{a},\ \ \ \ \ \ \ \ \ \ \left(  f^{a}\right)  ^{b}=f^{ab}.
\]

\end{theorem}

\begin{proof}
Easy exercise (Exercise \ref{exe.fps.power-c.rules}).
\end{proof}

Now, let us return to Example 2 from Section \ref{sec.gf.exas}. In that
example, we had to solve the quadratic equation
\[
C\left(  x\right)  =1+x\left(  C\left(  x\right)  \right)  ^{2}%
\ \ \ \ \ \ \ \ \ \ \text{for an FPS }C\left(  x\right)  \in\mathbb{Q}\left[
\left[  x\right]  \right]  .
\]
Let us write $C$ for $C\left(  x\right)  $; thus, this quadratic equation
becomes%
\[
C=1+xC^{2}.
\]
By completing the square, we can rewrite this equation in the equivalent form%
\[
\left(  1-2xC\right)  ^{2}=1-4x.
\]
Taking both sides of this equation to the $1/2$-th power, we obtain%
\[
\left(  \left(  1-2xC\right)  ^{2}\right)  ^{1/2}=\left(  1-4x\right)  ^{1/2}%
\]
(since both sides are FPSs with constant term $1$). However, the FPS $1-2xC$
has constant term $1$; thus, the rules of exponents yield $\left(  \left(
1-2xC\right)  ^{2}\right)  ^{1/2}=\left(  1-2xC\right)  ^{2\cdot1/2}=1-2xC$.
Hence,%
\[
1-2xC=\left(  \left(  1-2xC\right)  ^{2}\right)  ^{1/2}=\left(  1-4x\right)
^{1/2}.
\]
This is a linear equation in $C$; solving it for $C$ yields%
\[
C=\dfrac{1}{2x}\left(  1-\left(  1-4x\right)  ^{1/2}\right)  .
\]
This is precisely the \textquotedblleft square-root\textquotedblright%
\ expression for $C=C\left(  x\right)  $ that we have obtained back in Section
\ref{sec.gf.exas}, but now we have proved it rigorously.

\subsubsection{The Newton binomial formula for arbitrary exponents}

Is Example 2 from Section \ref{sec.gf.exas} fully justified now? No, because
we still need to prove the identity (\ref{eq.sec.gf.exas.2.(1+x)1/2}) that we
used back there. Since we are defining powers in the 2nd way (i.e., using
Definition \ref{def.fps.power-c} rather than using
(\ref{eq.sec.gf.nips.newton-def2})), it is not immediately obvious.
Nevertheless, it can be proved. More generally, we can prove the following:

\begin{theorem}
[Generalized Newton binomial formula]\label{thm.fps.gen-newton}Assume that $K$
is a commutative $\mathbb{Q}$-algebra. Let $c\in K$. Then,%
\[
\left(  1+x\right)  ^{c}=\sum_{k\in\mathbb{N}}\dbinom{c}{k}x^{k}.
\]

\end{theorem}

The following proof illustrates a technique that will probably appear
preposterous if you are seeing it for the first time, but is in fact both
legitimate and rather useful.

\begin{proof}
[Proof of Theorem \ref{thm.fps.gen-newton} (sketched).]The definition of
$\operatorname*{Log}$ yields
\[
\operatorname*{Log}\left(  1+x\right)  =\overline{\log}\circ\left(
\underbrace{\left(  1+x\right)  -1}_{=x}\right)  =\overline{\log}\circ
x=\overline{\log}%
\]
(by Proposition \ref{prop.fps.subs.rules} \textbf{(g)}, applied to
$g=\overline{\log}$).

Now, let us obstinately compute $\left(  1+x\right)  ^{c}$ using Definition
\ref{def.fps.power-c} and the definitions of $\operatorname*{Exp}$ and
$\operatorname*{Log}$. To wit: Let $\mathbb{P}$ denote the set $\left\{
1,2,3,\ldots\right\}  $. By Definition \ref{def.fps.power-c}, we have%
\begin{align}
&  \left(  1+x\right)  ^{c}\nonumber\\
&  =\operatorname*{Exp}\left(  c\operatorname*{Log}\left(  1+x\right)
\right)  =\operatorname*{Exp}\left(  c\,\overline{\log}\right)
\ \ \ \ \ \ \ \ \ \ \left(  \text{since }\operatorname*{Log}\left(
1+x\right)  =\overline{\log}\right) \nonumber\\
&  =\exp\circ\left(  c\,\overline{\log}\right)  \ \ \ \ \ \ \ \ \ \ \left(
\text{by the definition of }\operatorname*{Exp}\right) \nonumber\\
&  =\exp\circ\left(  c\sum_{n\geq1}\dfrac{\left(  -1\right)  ^{n-1}}{n}%
x^{n}\right)  \ \ \ \ \ \ \ \ \ \ \left(  \text{since }\overline{\log}%
=\sum_{n\geq1}\dfrac{\left(  -1\right)  ^{n-1}}{n}x^{n}\right) \nonumber\\
&  =\exp\circ\left(  \sum_{n\geq1}\dfrac{\left(  -1\right)  ^{n-1}}{n}%
cx^{n}\right) \nonumber\\
&  =\sum_{m\in\mathbb{N}}\dfrac{1}{m!}\left(  \sum_{n\geq1}\dfrac{\left(
-1\right)  ^{n-1}}{n}cx^{n}\right)  ^{m} \label{pf.thm.fps.gen-newton.1}%
\end{align}
(by Definition \ref{def.fps.subs}, since $\exp=\sum_{n\in\mathbb{N}}\dfrac
{1}{n!}x^{n}=\sum_{m\in\mathbb{N}}\dfrac{1}{m!}x^{m}$).

Now, fix $m\in\mathbb{N}$. We shall expand $\left(  \sum_{n\geq1}%
\dfrac{\left(  -1\right)  ^{n-1}}{n}cx^{n}\right)  ^{m}$. Indeed, we can
replace the \textquotedblleft$\sum_{n\geq1}$\textquotedblright\ sign by an
\textquotedblleft$\sum_{n\in\mathbb{P}}$\textquotedblright\ sign, since
$\mathbb{P}=\left\{  1,2,3,\ldots\right\}  $. Thus,%
\begin{align*}
&  \left(  \sum_{n\geq1}\dfrac{\left(  -1\right)  ^{n-1}}{n}cx^{n}\right)
^{m}\\
&  =\left(  \sum_{n\in\mathbb{P}}\dfrac{\left(  -1\right)  ^{n-1}}{n}%
cx^{n}\right)  ^{m}\\
&  =\underbrace{\left(  \sum_{n\in\mathbb{P}}\dfrac{\left(  -1\right)  ^{n-1}%
}{n}cx^{n}\right)  \left(  \sum_{n\in\mathbb{P}}\dfrac{\left(  -1\right)
^{n-1}}{n}cx^{n}\right)  \cdots\left(  \sum_{n\in\mathbb{P}}\dfrac{\left(
-1\right)  ^{n-1}}{n}cx^{n}\right)  }_{m\text{ times}}\\
&  =\left(  \sum_{n_{1}\in\mathbb{P}}\dfrac{\left(  -1\right)  ^{n_{1}-1}%
}{n_{1}}cx^{n_{1}}\right)  \left(  \sum_{n_{2}\in\mathbb{P}}\dfrac{\left(
-1\right)  ^{n_{2}-1}}{n_{2}}cx^{n_{2}}\right)  \cdots\left(  \sum_{n_{m}%
\in\mathbb{P}}\dfrac{\left(  -1\right)  ^{n_{m}-1}}{n_{m}}cx^{n_{m}}\right) \\
&  \ \ \ \ \ \ \ \ \ \ \ \ \ \ \ \ \ \ \ \ \left(  \text{here, we have renamed
the summation indices}\right) \\
&  =\sum_{\left(  n_{1},n_{2},\ldots,n_{m}\right)  \in\mathbb{P}^{m}}\left(
\dfrac{\left(  -1\right)  ^{n_{1}-1}}{n_{1}}cx^{n_{1}}\right)  \left(
\dfrac{\left(  -1\right)  ^{n_{2}-1}}{n_{2}}cx^{n_{2}}\right)  \cdots\left(
\dfrac{\left(  -1\right)  ^{n_{m}-1}}{n_{m}}cx^{n_{m}}\right)
\end{align*}
(by a product rule for the product of $m$ sums\footnote{This product rule says
that
\[
\left(  \sum_{n_{1}\in A_{1}}a_{1,n_{1}}\right)  \left(  \sum_{n_{2}\in A_{2}%
}a_{2,n_{2}}\right)  \cdots\left(  \sum_{n_{m}\in A_{m}}a_{m,n_{m}}\right)
=\sum_{\left(  n_{1},n_{2},\ldots,n_{m}\right)  \in A_{1}\times A_{2}%
\times\cdots\times A_{m}}a_{1,n_{1}}a_{2,n_{2}}\cdots a_{m,n_{m}}%
\]
for any $m$ sets $A_{1},A_{2},\ldots,A_{m}$ and any elements $a_{i,j}\in K$,
provided that all the sums on the left hand side of this equality are
summable. We leave it to the reader to convince himself of this rule
(intuitively, it just says that we can expand a product of sums in the usual
way, even when the sums are infinite) and to check that the sums we are
applying it to are indeed summable.}). Hence,%
\begin{align}
&  \left(  \sum_{n\geq1}\dfrac{\left(  -1\right)  ^{n-1}}{n}cx^{n}\right)
^{m}\nonumber\\
&  =\underbrace{\sum_{\left(  n_{1},n_{2},\ldots,n_{m}\right)  \in
\mathbb{P}^{m}}}_{=\sum_{k\in\mathbb{N}}\ \ \sum_{\substack{\left(
n_{1},n_{2},\ldots,n_{m}\right)  \in\mathbb{P}^{m};\\n_{1}+n_{2}+\cdots
+n_{m}=k}}}\underbrace{\left(  \dfrac{\left(  -1\right)  ^{n_{1}-1}}{n_{1}%
}cx^{n_{1}}\right)  \left(  \dfrac{\left(  -1\right)  ^{n_{2}-1}}{n_{2}%
}cx^{n_{2}}\right)  \cdots\left(  \dfrac{\left(  -1\right)  ^{n_{m}-1}}{n_{m}%
}cx^{n_{m}}\right)  }_{=\dfrac{\left(  -1\right)  ^{n_{1}+n_{2}+\cdots
+n_{m}-m}}{n_{1}n_{2}\cdots n_{m}}c^{m}x^{n_{1}+n_{2}+\cdots+n_{m}}%
}\nonumber\\
&  =\sum_{k\in\mathbb{N}}\ \ \sum_{\substack{\left(  n_{1},n_{2},\ldots
,n_{m}\right)  \in\mathbb{P}^{m};\\n_{1}+n_{2}+\cdots+n_{m}=k}}\dfrac{\left(
-1\right)  ^{n_{1}+n_{2}+\cdots+n_{m}-m}}{n_{1}n_{2}\cdots n_{m}}%
c^{m}\underbrace{x^{n_{1}+n_{2}+\cdots+n_{m}}}_{\substack{=x^{k}\\\text{(since
}n_{1}+n_{2}+\cdots+n_{m}=k\text{)}}}\nonumber\\
&  =\sum_{k\in\mathbb{N}}\ \ \sum_{\substack{\left(  n_{1},n_{2},\ldots
,n_{m}\right)  \in\mathbb{P}^{m};\\n_{1}+n_{2}+\cdots+n_{m}=k}}\dfrac{\left(
-1\right)  ^{n_{1}+n_{2}+\cdots+n_{m}-m}}{n_{1}n_{2}\cdots n_{m}}c^{m}x^{k}.
\label{pf.thm.fps.gen-newton.2}%
\end{align}


Now, forget that we fixed $m$. We thus have proved
(\ref{pf.thm.fps.gen-newton.2}) for each $m\in\mathbb{N}$.

Now, (\ref{pf.thm.fps.gen-newton.1}) becomes%
\begin{align}
&  \left(  1+x\right)  ^{c}\nonumber\\
&  =\sum_{m\in\mathbb{N}}\dfrac{1}{m!}\left(  \sum_{n\geq1}\dfrac{\left(
-1\right)  ^{n-1}}{n}cx^{n}\right)  ^{m}\nonumber\\
&  =\sum_{m\in\mathbb{N}}\dfrac{1}{m!}\sum_{k\in\mathbb{N}}\ \ \sum
_{\substack{\left(  n_{1},n_{2},\ldots,n_{m}\right)  \in\mathbb{P}^{m}%
;\\n_{1}+n_{2}+\cdots+n_{m}=k}}\dfrac{\left(  -1\right)  ^{n_{1}+n_{2}%
+\cdots+n_{m}-m}}{n_{1}n_{2}\cdots n_{m}}c^{m}x^{k}\ \ \ \ \ \ \ \ \ \ \left(
\text{by (\ref{pf.thm.fps.gen-newton.2})}\right) \nonumber\\
&  =\sum_{k\in\mathbb{N}}\ \ \sum_{m\in\mathbb{N}}\ \ \sum_{\substack{\left(
n_{1},n_{2},\ldots,n_{m}\right)  \in\mathbb{P}^{m};\\n_{1}+n_{2}+\cdots
+n_{m}=k}}\dfrac{1}{m!}\cdot\dfrac{\left(  -1\right)  ^{n_{1}+n_{2}%
+\cdots+n_{m}-m}}{n_{1}n_{2}\cdots n_{m}}c^{m}x^{k}\nonumber\\
&  =\sum_{k\in\mathbb{N}}\left(  \sum_{m\in\mathbb{N}}\ \ \sum
_{\substack{\left(  n_{1},n_{2},\ldots,n_{m}\right)  \in\mathbb{P}^{m}%
;\\n_{1}+n_{2}+\cdots+n_{m}=k}}\dfrac{1}{m!}\cdot\dfrac{\left(  -1\right)
^{n_{1}+n_{2}+\cdots+n_{m}-m}}{n_{1}n_{2}\cdots n_{m}}c^{m}\right)  x^{k}.
\label{pf.thm.fps.gen-newton.3}%
\end{align}


Now, let $k\in\mathbb{N}$. Let us rewrite the \textquotedblleft middle
sum\textquotedblright\ $\sum_{m\in\mathbb{N}}\ \ \sum_{\substack{\left(
n_{1},n_{2},\ldots,n_{m}\right)  \in\mathbb{P}^{m};\\n_{1}+n_{2}+\cdots
+n_{m}=k}}\dfrac{1}{m!}\cdot\dfrac{\left(  -1\right)  ^{n_{1}+n_{2}%
+\cdots+n_{m}-m}}{n_{1}n_{2}\cdots n_{m}}c^{m}$ on the right hand side as a
finite sum. Indeed, a \emph{composition} of $k$ shall mean a tuple $\left(
n_{1},n_{2},\ldots,n_{m}\right)  $ of positive integers satisfying
$n_{1}+n_{2}+\cdots+n_{m}=k$. (For example, $\left(  1,3,1\right)  $ is a
composition of $5$. We will study compositions in more detail in Section
\ref{sec.fps.intcomps}.) Let $\operatorname*{Comp}\left(  k\right)  $ denote
the set of all compositions of $k$. It is easy to see that this set
$\operatorname*{Comp}\left(  k\right)  $ is finite\footnote{\textit{Proof.}
Let $\left(  n_{1},n_{2},\ldots,n_{m}\right)  \in\operatorname*{Comp}\left(
k\right)  $. Thus, $\left(  n_{1},n_{2},\ldots,n_{m}\right)  $ is a
composition of $k$. In other words, $\left(  n_{1},n_{2},\ldots,n_{m}\right)
$ is a finite tuple of positive integers satisfying $n_{1}+n_{2}+\cdots
+n_{m}=k$. Hence, all its $m$ entries $n_{1},n_{2},\ldots,n_{m}$ are positive
integers and thus are $\geq1$; therefore, $n_{1}+n_{2}+\cdots+n_{m}%
\geq\underbrace{1+1+\cdots+1}_{m\text{ times}}=m$, so that $m\leq n_{1}%
+n_{2}+\cdots+n_{m}=k$. Thus, $m\in\left\{  0,1,\ldots,k\right\}  $.
\par
Furthermore, the sum $n_{1}+n_{2}+\cdots+n_{m}$ is $\geq$ to each of its $m$
addends (since its $m$ addends $n_{1},n_{2},\ldots,n_{m}$ are positive). In
other words, we have $n_{1}+n_{2}+\cdots+n_{m}\geq n_{i}$ for each
$i\in\left\{  1,2,\ldots,m\right\}  $. Thus, for each $i\in\left\{
1,2,\ldots,m\right\}  $, we have $n_{i}\leq n_{1}+n_{2}+\cdots+n_{m}=k$ and
therefore $n_{i}\in\left\{  1,2,\ldots,k\right\}  $ (since $n_{i}$ is a
positive integer). Hence,
\[
\left(  n_{1},n_{2},\ldots,n_{m}\right)  \in\left\{  1,2,\ldots,k\right\}
^{m}\subseteq\bigcup_{\ell\in\left\{  0,1,\ldots,k\right\}  }\left\{
1,2,\ldots,k\right\}  ^{\ell}%
\]
(since $m\in\left\{  0,1,\ldots,k\right\}  $).
\par
Now, forget that we fixed $\left(  n_{1},n_{2},\ldots,n_{m}\right)  $. We thus
have shown that $\left(  n_{1},n_{2},\ldots,n_{m}\right)  \in\bigcup_{\ell
\in\left\{  0,1,\ldots,k\right\}  }\left\{  1,2,\ldots,k\right\}  ^{\ell}$ for
each $\left(  n_{1},n_{2},\ldots,n_{m}\right)  \in\operatorname*{Comp}\left(
k\right)  $. In other words, $\operatorname*{Comp}\left(  k\right)
\subseteq\bigcup_{\ell\in\left\{  0,1,\ldots,k\right\}  }\left\{
1,2,\ldots,k\right\}  ^{\ell}$. Since the set $\bigcup_{\ell\in\left\{
0,1,\ldots,k\right\}  }\left\{  1,2,\ldots,k\right\}  ^{\ell}$ is clearly
finite (having size $\sum_{\ell\in\left\{  0,1,\ldots,k\right\}  }k^{\ell}$),
this entails that the set $\operatorname*{Comp}\left(  k\right)  $ is finite
as well, qed.
\par
(Incidentally, we will see in Section \ref{sec.fps.intcomps} that this set
$\operatorname*{Comp}\left(  k\right)  $ has size $2^{k-1}$ for $k\geq1$, and
size $1$ for $k=0$.)}. Now, we can rewrite the double summation sign
\textquotedblleft$\sum_{m\in\mathbb{N}}\ \ \sum_{\substack{\left(  n_{1}%
,n_{2},\ldots,n_{m}\right)  \in\mathbb{P}^{m};\\n_{1}+n_{2}+\cdots+n_{m}=k}%
}$\textquotedblright\ as a single summation sign \textquotedblleft%
$\sum_{\left(  n_{1},n_{2},\ldots,n_{m}\right)  \in\operatorname*{Comp}\left(
k\right)  }$\textquotedblright\ (since $\operatorname*{Comp}\left(  k\right)
$ is precisely the set of all tuples $\left(  n_{1},n_{2},\ldots,n_{m}\right)
\in\mathbb{P}^{m}$ satisfying $n_{1}+n_{2}+\cdots+n_{m}=k$). Hence, we obtain
\begin{align}
&  \sum_{m\in\mathbb{N}}\ \ \sum_{\substack{\left(  n_{1},n_{2},\ldots
,n_{m}\right)  \in\mathbb{P}^{m};\\n_{1}+n_{2}+\cdots+n_{m}=k}}\dfrac{1}%
{m!}\cdot\dfrac{\left(  -1\right)  ^{n_{1}+n_{2}+\cdots+n_{m}-m}}{n_{1}%
n_{2}\cdots n_{m}}c^{m}\nonumber\\
&  =\sum_{\left(  n_{1},n_{2},\ldots,n_{m}\right)  \in\operatorname*{Comp}%
\left(  k\right)  }\dfrac{1}{m!}\cdot\dfrac{\left(  -1\right)  ^{n_{1}%
+n_{2}+\cdots+n_{m}-m}}{n_{1}n_{2}\cdots n_{m}}c^{m}.
\label{pf.thm.fps.gen-newton.4}%
\end{align}


Forget that we fixed $k$. Thus, for each $k\in\mathbb{N}$, we have defined a
finite set $\operatorname*{Comp}\left(  k\right)  $ and shown that
(\ref{pf.thm.fps.gen-newton.4}) holds.

Using (\ref{pf.thm.fps.gen-newton.4}), we can rewrite
(\ref{pf.thm.fps.gen-newton.3}) as%
\begin{align}
&  \left(  1+x\right)  ^{c}\nonumber\\
&  =\sum_{k\in\mathbb{N}}\left(  \sum_{\left(  n_{1},n_{2},\ldots
,n_{m}\right)  \in\operatorname*{Comp}\left(  k\right)  }\dfrac{1}{m!}%
\cdot\dfrac{\left(  -1\right)  ^{n_{1}+n_{2}+\cdots+n_{m}-m}}{n_{1}n_{2}\cdots
n_{m}}c^{m}\right)  x^{k}. \label{pf.thm.fps.gen-newton.5}%
\end{align}


Now, recall that our goal is to prove that this equals
\[
\sum_{k\in\mathbb{N}}\dbinom{c}{k}x^{k}.
\]
This is equivalent to proving that the equality%
\begin{equation}
\sum_{\left(  n_{1},n_{2},\ldots,n_{m}\right)  \in\operatorname*{Comp}\left(
k\right)  }\dfrac{1}{m!}\cdot\dfrac{\left(  -1\right)  ^{n_{1}+n_{2}%
+\cdots+n_{m}-m}}{n_{1}n_{2}\cdots n_{m}}c^{m}=\dbinom{c}{k}
\label{pf.thm.fps.gen-newton.polyeq}%
\end{equation}
holds for each $k\in\mathbb{N}$ (because two FPSs $u=\sum_{k\in\mathbb{N}%
}u_{k}x^{k}$ and $v=\sum_{k\in\mathbb{N}}v_{k}x^{k}$ (with $u_{k}\in K$ and
$v_{k}\in K$) are equal if and only if the equality $u_{k}=v_{k}$ holds for
each $k\in\mathbb{N}$).

Thus, we have reduced our original goal (which was to prove $\left(
1+x\right)  ^{c}=\sum_{k\in\mathbb{N}}\dbinom{c}{k}x^{k}$) to the auxiliary
goal of proving the equality (\ref{pf.thm.fps.gen-newton.polyeq}) for each
$k\in\mathbb{N}$. However, this doesn't look very useful, since
(\ref{pf.thm.fps.gen-newton.polyeq}) is too messy an equality to have a simple
proof. We are seemingly stuck.

However, it turns out that we are almost there -- we just need to take a
bird's eye view. Here is the plan: We fix $k\in\mathbb{N}$. Instead of trying
to prove the equality (\ref{pf.thm.fps.gen-newton.polyeq}) directly, we
observe that both sides of this equality are polynomials (with rational
coefficients) in $c$. (Indeed, the left hand side is clearly a polynomial in
$c$, since it is a finite sum of \textquotedblleft rational number times a
power of $c$\textquotedblright\ expressions. The right hand side is a
polynomial in $c$ because $\dbinom{c}{k}=\dfrac{c\left(  c-1\right)  \left(
c-2\right)  \cdots\left(  c-k+1\right)  }{k!}$.) Thus, the polynomial identity
trick (which we learnt in Subsection \ref{subsec.gf.defs.cvi}) tells us that
if we can prove this equality (\ref{pf.thm.fps.gen-newton.polyeq}) for each
$c\in\mathbb{N}$, then it will automatically hold for each $c\in K$ (since the
two polynomials that yield its left and right hand sides will have to be
equal, having infinitely many equal values). Hence, in order to prove
(\ref{pf.thm.fps.gen-newton.polyeq}) for each $c\in K$, it suffices to prove
it for each $c\in\mathbb{N}$. Now, how can we prove it for each $c\in
\mathbb{N}$ ? We forget that we fixed $k$, and we remember that the equality
(\ref{pf.thm.fps.gen-newton.polyeq}) (for all $k\in\mathbb{N}$) is just an
equivalent restatement of the FPS equality $\left(  1+x\right)  ^{c}%
=\sum_{k\in\mathbb{N}}\dbinom{c}{k}x^{k}$ (that is, the equality we have
originally set out to prove). However, we know for sure that this equality
holds for each $c\in\mathbb{N}$ (by Theorem \ref{thm.fps.newton-binom},
applied to $n=c$). Hence, the equality (\ref{pf.thm.fps.gen-newton.polyeq})
also holds for each $c\in\mathbb{N}$ (and each $k\in\mathbb{N}$). And this is
precisely what we needed to show!

Let me explain this argument in detail now, as it is somewhat
vertigo-inducing. We forget that we fixed $K$ and $c$. Now, fix $c\in
\mathbb{N}$. Thus, $c\in\mathbb{N}\subseteq\mathbb{Z}$. Hence, in the ring
$\mathbb{Q}\left[  \left[  x\right]  \right]  $, we have%
\[
\left(  1+x\right)  ^{c}=\sum_{k\in\mathbb{N}}\dbinom{c}{k}x^{k}%
\ \ \ \ \ \ \ \ \ \ \left(  \text{by Theorem \ref{thm.fps.newton-binom},
applied to }n=c\right)  .
\]
However, (\ref{pf.thm.fps.gen-newton.5}) (applied to $K=\mathbb{Q}$) shows
that%
\[
\left(  1+x\right)  ^{c}=\sum_{k\in\mathbb{N}}\left(  \sum_{\left(
n_{1},n_{2},\ldots,n_{m}\right)  \in\operatorname*{Comp}\left(  k\right)
}\dfrac{1}{m!}\cdot\dfrac{\left(  -1\right)  ^{n_{1}+n_{2}+\cdots+n_{m}-m}%
}{n_{1}n_{2}\cdots n_{m}}c^{m}\right)  x^{k}.
\]
Comparing these two equalities, we obtain%
\[
\sum_{k\in\mathbb{N}}\left(  \sum_{\left(  n_{1},n_{2},\ldots,n_{m}\right)
\in\operatorname*{Comp}\left(  k\right)  }\dfrac{1}{m!}\cdot\dfrac{\left(
-1\right)  ^{n_{1}+n_{2}+\cdots+n_{m}-m}}{n_{1}n_{2}\cdots n_{m}}c^{m}\right)
x^{k}=\sum_{k\in\mathbb{N}}\dbinom{c}{k}x^{k}.
\]
Comparing coefficients in this equality, we see that%
\begin{equation}
\sum_{\left(  n_{1},n_{2},\ldots,n_{m}\right)  \in\operatorname*{Comp}\left(
k\right)  }\dfrac{1}{m!}\cdot\dfrac{\left(  -1\right)  ^{n_{1}+n_{2}%
+\cdots+n_{m}-m}}{n_{1}n_{2}\cdots n_{m}}c^{m}=\dbinom{c}{k}
\label{pf.thm.fps.gen-newton.finale.1}%
\end{equation}
for each $k\in\mathbb{N}$. This is an equality between two rational numbers.

Now, forget that we fixed $c$. We thus have shown that
(\ref{pf.thm.fps.gen-newton.finale.1}) holds for each $k\in\mathbb{N}$ and
each $c\in\mathbb{N}$.

Let us now fix $k\in\mathbb{N}$. We have just shown that the equality
(\ref{pf.thm.fps.gen-newton.finale.1}) holds for each $c\in\mathbb{N}$. In
other words, the two polynomials%
\[
f:=\sum_{\left(  n_{1},n_{2},\ldots,n_{m}\right)  \in\operatorname*{Comp}%
\left(  k\right)  }\dfrac{1}{m!}\cdot\dfrac{\left(  -1\right)  ^{n_{1}%
+n_{2}+\cdots+n_{m}-m}}{n_{1}n_{2}\cdots n_{m}}x^{m}\in\mathbb{Q}\left[
x\right]
\]
and%
\[
g:=\dbinom{x}{k}=\dfrac{x\left(  x-1\right)  \left(  x-2\right)  \cdots\left(
x-k+1\right)  }{k!}\in\mathbb{Q}\left[  x\right]
\]
satisfy $f\left[  c\right]  =g\left[  c\right]  $ for each $c\in\mathbb{N}$
(because $f\left[  c\right]  $ is the left hand side of
(\ref{pf.thm.fps.gen-newton.finale.1}), while $g\left[  c\right]  $ is the
right hand side of (\ref{pf.thm.fps.gen-newton.finale.1})). Thus, each
$c\in\mathbb{N}$ satisfies $\left(  f-g\right)  \left[  c\right]
=\underbrace{f\left[  c\right]  }_{=g\left[  c\right]  }-\,g\left[  c\right]
=g\left[  c\right]  -g\left[  c\right]  =0$. In other words, each
$c\in\mathbb{N}$ is a root of $f-g$. Hence, the polynomial $f-g$ has
infinitely many roots in $\mathbb{Q}$ (since there are infinitely many
$c\in\mathbb{N}$). Since $f-g$ is a polynomial with rational coefficients,
this is impossible unless $f-g=0$. We thus must have $f-g=0$, so that $f=g$.
In other words,
\begin{equation}
\sum_{\left(  n_{1},n_{2},\ldots,n_{m}\right)  \in\operatorname*{Comp}\left(
k\right)  }\dfrac{1}{m!}\cdot\dfrac{\left(  -1\right)  ^{n_{1}+n_{2}%
+\cdots+n_{m}-m}}{n_{1}n_{2}\cdots n_{m}}x^{m}=\dbinom{x}{k}
\label{pf.thm.fps.gen-newton.finale.4}%
\end{equation}
holds in the polynomial ring $\mathbb{Q}\left[  x\right]  $.

Now, forget that we fixed $k$. We thus have shown that the equality
(\ref{pf.thm.fps.gen-newton.finale.4}) holds for each $k\in\mathbb{N}$.

Now, fix a commutative $\mathbb{Q}$-algebra $K$ and an arbitrary element $c\in
K$. For each $k\in\mathbb{N}$, we then have%
\begin{equation}
\sum_{\left(  n_{1},n_{2},\ldots,n_{m}\right)  \in\operatorname*{Comp}\left(
k\right)  }\dfrac{1}{m!}\cdot\dfrac{\left(  -1\right)  ^{n_{1}+n_{2}%
+\cdots+n_{m}-m}}{n_{1}n_{2}\cdots n_{m}}c^{m}=\dbinom{c}{k}
\label{pf.thm.fps.gen-newton.finale.5}%
\end{equation}
(by substituting $c$ for $x$ on both sides of the equality
(\ref{pf.thm.fps.gen-newton.finale.4})). Consequently, we can rewrite
(\ref{pf.thm.fps.gen-newton.5}) as%
\[
\left(  1+x\right)  ^{c}=\sum_{k\in\mathbb{N}}\dbinom{c}{k}x^{k}.
\]
This proves Theorem \ref{thm.fps.gen-newton}.
\end{proof}

The method we used in the above proof is worth recapitulating in broad strokes:

\begin{itemize}
\item We had to prove a fairly abstract statement (namely, the identity
$\left(  1+x\right)  ^{c}=\sum_{k\in\mathbb{N}}\dbinom{c}{k}x^{k}$).

\item We translated this statement into an awkward but more concrete statement
(namely, the equality (\ref{pf.thm.fps.gen-newton.polyeq})).

\item We then argued that this concrete statement needs only to be proven in a
special case (viz., for all $c\in\mathbb{N}$ rather than for all $c\in K$),
because it is an equality between two polynomials with rational coefficients.

\item To prove this concrete statement in this special case, we translated it
back into the abstract language of FPSs, and realized that in this special
case it is already known (as a consequence of Theorem
\ref{thm.fps.newton-binom}).
\end{itemize}

Thus, by strategically switching between the abstract and the concrete, we
have managed to use the advantages of both sides.

Now that Theorem \ref{thm.fps.gen-newton} is proved, Example 2 from Section
\ref{sec.gf.exas} is fully justified (since we can obtain
(\ref{eq.sec.gf.exas.2.(1+x)1/2}) by applying Theorem \ref{thm.fps.gen-newton}
to $K=\mathbb{Q}$ and $c=1/2$).

\subsubsection{Another application}

Let us show yet another application of powers with non-integer exponents and
the generalized Newton formula. We shall show the following binomial identity:

\begin{proposition}
\label{prop.binom.nCk-2i-qedmo.CN}Let $n\in\mathbb{C}$ and $k\in\mathbb{N}$.
Then,%
\[
\sum_{i=0}^{k}\dbinom{n+i-1}{i}\dbinom{n}{k-2i}=\dbinom{n+k-1}{k}.
\]

\end{proposition}

Proposition \ref{prop.binom.nCk-2i-qedmo.CN} can be proved in various ways.
For example, a mostly combinatorial proof is found in \cite[Exercise 2.10.7
and Exercise 2.10.8]{19fco}\footnote{Specifically, \cite[Exercise
2.10.7]{19fco} proves Proposition \ref{prop.binom.nCk-2i-qedmo.CN} in the
particular case when $n\in\mathbb{N}$; then, \cite[Exercise 2.10.8]{19fco}
extends it to the case when $n\in\mathbb{R}$. However, the latter argument can
just as well be used to extend it to arbitrary $n\in\mathbb{C}$.}. We shall
give a proof using generating functions instead.

\begin{proof}
[Proof of Proposition \ref{prop.binom.nCk-2i-qedmo.CN}.]Define two FPSs
$f,g\in\mathbb{C}\left[  \left[  x\right]  \right]  $ by
\begin{align}
f  &  =\sum_{i\in\mathbb{N}}\dbinom{n+i-1}{i}x^{2i}%
\label{pf.prop.binom.nCk-2i-qedmo.CN.f=}\\
\text{and}\ \ \ \ \ \ \ \ \ \ g  &  =\sum_{j\in\mathbb{N}}\dbinom{n}{j}x^{j}.
\label{pf.prop.binom.nCk-2i-qedmo.CN.g=}%
\end{align}
(We will soon see why we chose to define them this way.) Multiplying the two
equalities (\ref{pf.prop.binom.nCk-2i-qedmo.CN.f=}) and
(\ref{pf.prop.binom.nCk-2i-qedmo.CN.g=}), we find%
\begin{align*}
fg  &  =\left(  \sum_{i\in\mathbb{N}}\dbinom{n+i-1}{i}x^{2i}\right)  \left(
\sum_{j\in\mathbb{N}}\dbinom{n}{j}x^{j}\right) \\
&  =\underbrace{\sum_{i\in\mathbb{N}}\ \ \sum_{j\in\mathbb{N}}}_{=\sum
_{\left(  i,j\right)  \in\mathbb{N}\times\mathbb{N}}}\ \ \underbrace{\dbinom
{n+i-1}{i}x^{2i}\dbinom{n}{j}x^{j}}_{=\dbinom{n+i-1}{i}\dbinom{n}{j}x^{2i+j}%
}\\
&  =\underbrace{\sum_{\left(  i,j\right)  \in\mathbb{N}\times\mathbb{N}}%
}_{=\sum_{m\in\mathbb{N}}\ \ \sum_{\substack{\left(  i,j\right)  \in
\mathbb{N}\times\mathbb{N};\\2i+j=m}}}\dbinom{n+i-1}{i}\dbinom{n}{j}x^{2i+j}\\
&  =\sum_{m\in\mathbb{N}}\ \ \sum_{\substack{\left(  i,j\right)  \in
\mathbb{N}\times\mathbb{N};\\2i+j=m}}\dbinom{n+i-1}{i}\dbinom{n}%
{j}\underbrace{x^{2i+j}}_{\substack{=x^{m}\\\text{(since }2i+j=m\text{)}}}\\
&  =\sum_{m\in\mathbb{N}}\ \ \sum_{\substack{\left(  i,j\right)  \in
\mathbb{N}\times\mathbb{N};\\2i+j=m}}\dbinom{n+i-1}{i}\dbinom{n}{j}x^{m}\\
&  =\sum_{m\in\mathbb{N}}\left(  \sum_{\substack{\left(  i,j\right)
\in\mathbb{N}\times\mathbb{N};\\2i+j=m}}\dbinom{n+i-1}{i}\dbinom{n}{j}\right)
x^{m}.
\end{align*}
Hence, the $x^{k}$-coefficient of this FPS $fg$ is%
\begin{equation}
\left[  x^{k}\right]  \left(  fg\right)  =\sum_{\substack{\left(  i,j\right)
\in\mathbb{N}\times\mathbb{N};\\2i+j=k}}\dbinom{n+i-1}{i}\dbinom{n}{j}.
\label{pf.prop.binom.nCk-2i-qedmo.CN.3}%
\end{equation}


Now, a pair $\left(  i,j\right)  \in\mathbb{N}\times\mathbb{N}$ satisfying
$2i+j=k$ is uniquely determined by its first entry $i$, since its second entry
$j$ is given by $j=k-2i$. Hence, we can substitute $\left(  i,k-2i\right)  $
for $\left(  i,j\right)  $ in the sum $\sum_{\substack{\left(  i,j\right)
\in\mathbb{N}\times\mathbb{N};\\2i+j=k}}\dbinom{n+i-1}{i}\dbinom{n}{j}$, thus
rewriting this sum as $\sum_{\substack{i\in\mathbb{N};\\k-2i\in\mathbb{N}%
}}\dbinom{n+i-1}{i}\dbinom{n}{k-2i}$. Hence,%
\begin{align*}
\sum_{\substack{\left(  i,j\right)  \in\mathbb{N}\times\mathbb{N}%
;\\2i+j=k}}\dbinom{n+i-1}{i}\dbinom{n}{j}  &  =\sum_{\substack{i\in
\mathbb{N};\\k-2i\in\mathbb{N}}}\dbinom{n+i-1}{i}\dbinom{n}{k-2i}\\
&  =\sum_{\substack{i\in\mathbb{N};\\2i\leq k}}\dbinom{n+i-1}{i}\dbinom
{n}{k-2i}\\
&  \ \ \ \ \ \ \ \ \ \ \ \ \ \ \ \ \ \ \ \ \left(
\begin{array}
[c]{c}%
\text{since an }i\in\mathbb{N}\text{ satisfies }k-2i\in\mathbb{N}\\
\text{if and only if it satisfies }2i\leq k
\end{array}
\right) \\
&  =\sum_{\substack{i\in\mathbb{N};\\i\leq k}}\dbinom{n+i-1}{i}\dbinom
{n}{k-2i}%
\end{align*}
(here, we have replaced the condition \textquotedblleft$2i\leq k$%
\textquotedblright\ under the summation sign by the weaker condition
\textquotedblleft$i\leq k$\textquotedblright, thus extending the range of the
sum; but this did not change the sum, since all the newly introduced addends
are $0$ because of the vanishing $\dbinom{n}{k-2i}$ factor). Thus,
(\ref{pf.prop.binom.nCk-2i-qedmo.CN.3}) becomes%
\begin{align}
\left[  x^{k}\right]  \left(  fg\right)   &  =\sum_{\substack{\left(
i,j\right)  \in\mathbb{N}\times\mathbb{N};\\2i+j=k}}\dbinom{n+i-1}{i}%
\dbinom{n}{j}\nonumber\\
&  =\sum_{\substack{i\in\mathbb{N};\\i\leq k}}\dbinom{n+i-1}{i}\dbinom
{n}{k-2i}\nonumber\\
&  =\sum_{i=0}^{k}\dbinom{n+i-1}{i}\dbinom{n}{k-2i}.
\label{pf.prop.binom.nCk-2i-qedmo.CN.5}%
\end{align}
Note that the right hand side here is precisely the left hand side of the
identity we are trying to prove. This is why we defined $f$ and $g$ as we did.
With a bit of experience, the computation above can easily be
reverse-engineered, and the definitions of $f$ and $g$ are essentially forced
by the goal of making (\ref{pf.prop.binom.nCk-2i-qedmo.CN.5}) hold.

Anyway, it is now clear that a simple expression for $fg$ would move us
forward. So let us try to simplify $f$ and $g$. For $g$, the answer is
easiest: We have%
\[
g=\sum_{j\in\mathbb{N}}\dbinom{n}{j}x^{j}=\left(  1+x\right)  ^{n},
\]
because Theorem \ref{thm.fps.gen-newton} (applied to $c=n$) yields $\left(
1+x\right)  ^{n}=\sum_{j\in\mathbb{N}}\dbinom{n}{j}x^{j}$. For $f$, we need a
few more steps. Proposition \ref{prop.fps.anti-newton-binom} yields%
\begin{equation}
\left(  1+x\right)  ^{-n}=\sum_{i\in\mathbb{N}}\left(  -1\right)  ^{i}%
\dbinom{n+i-1}{i}x^{i}. \label{pf.prop.binom.nCk-2i-qedmo.CN.9}%
\end{equation}
Substituting $-x^{2}$ for $x$ on both sides of this equality, we obtain%
\begin{align*}
\left(  1-x^{2}\right)  ^{-n}  &  =\sum_{i\in\mathbb{N}}\left(  -1\right)
^{i}\dbinom{n+i-1}{i}\underbrace{\left(  -x^{2}\right)  ^{i}}_{=\left(
-1\right)  ^{i}x^{2i}}=\sum_{i\in\mathbb{N}}\left(  -1\right)  ^{i}%
\dbinom{n+i-1}{i}\left(  -1\right)  ^{i}x^{2i}\\
&  =\sum_{i\in\mathbb{N}}\underbrace{\left(  -1\right)  ^{i}\left(  -1\right)
^{i}}_{=1}\dbinom{n+i-1}{i}x^{2i}=\sum_{i\in\mathbb{N}}\dbinom{n+i-1}{i}%
x^{2i}=f
\end{align*}
and thus $f=\left(  1-x^{2}\right)  ^{-n}$. Multiplying this equality by
$g=\left(  1+x\right)  ^{n}$, we obtain%
\begin{align*}
fg  &  =\left(  1-x^{2}\right)  ^{-n}\left(  1+x\right)  ^{n}=\left(
\dfrac{1+x}{1-x^{2}}\right)  ^{n}=\left(  \dfrac{1-x^{2}}{1+x}\right)  ^{-n}\\
&  =\left(  1-x\right)  ^{-n}\ \ \ \ \ \ \ \ \ \ \left(  \text{since }%
\dfrac{1-x^{2}}{1+x}=\dfrac{\left(  1-x\right)  \left(  1+x\right)  }%
{1+x}=1-x\right) \\
&  =\sum_{i\in\mathbb{N}}\left(  -1\right)  ^{i}\dbinom{n+i-1}{i}%
\underbrace{\left(  -x\right)  ^{i}}_{=\left(  -1\right)  ^{i}x^{i}}\\
&  \ \ \ \ \ \ \ \ \ \ \ \ \ \ \ \ \ \ \ \ \left(  \text{this follows by
substituting }-x\text{ for }x\text{ on both sides of
(\ref{pf.prop.binom.nCk-2i-qedmo.CN.9})}\right) \\
&  =\sum_{i\in\mathbb{N}}\left(  -1\right)  ^{i}\dbinom{n+i-1}{i}\left(
-1\right)  ^{i}x^{i}=\sum_{i\in\mathbb{N}}\underbrace{\left(  -1\right)
^{i}\left(  -1\right)  ^{i}}_{=1}\dbinom{n+i-1}{i}x^{i}\\
&  =\sum_{i\in\mathbb{N}}\dbinom{n+i-1}{i}x^{i}.
\end{align*}
Hence, the $x^{k}$-coefficient in $fg$ is $\left[  x^{k}\right]  \left(
fg\right)  =\dbinom{n+k-1}{k}$. Comparing this with
(\ref{pf.prop.binom.nCk-2i-qedmo.CN.5}), we obtain%
\[
\sum_{i=0}^{k}\dbinom{n+i-1}{i}\dbinom{n}{k-2i}=\dbinom{n+k-1}{k}.
\]
Thus, Proposition \ref{prop.binom.nCk-2i-qedmo.CN} is proved.
\end{proof}
